\section{Herramientas y Entorno de Desarrollo}

El desarrollo del proyecto se ha realizado principalmente utilizando el lenguaje de programación Python, empleando un entorno virtual gestionado mediante Conda y aceleración por GPU mediante CUDA para el entrenamiento de los modelos desarrollados con la biblioteca PyTorch.

El proceso de trabajo ha sido eminentemente iterativo y experimental, sin seguir una metodología formal concreta. En lugar de partir de un diseño definitivo, el proyecto avanzó mediante pruebas sucesivas documentadas en cuadernos de Jupyter Notebook. Estos cuadernos permiten registrar las principales etapas del desarrollo: exploración inicial del conjunto de datos, pruebas de diferentes representaciones del estado, construcción progresiva de la arquitectura Transformer, validación mediante pruebas de sobreajuste, entrenamiento a diferentes escalas y análisis estratégico de los resultados obtenidos.

Este enfoque ha facilitado la experimentación continua con distintos tamaños de conjunto de datos, configuraciones de entrenamiento y métricas de evaluación. Algunas decisiones importantes, como el uso de una representación mediante tokens o la incorporación de la evaluación estratégica, surgieron como resultado de estas iteraciones y no de una planificación cerrada desde el inicio.

Todo el código fuente del proyecto, así como los distintos experimentos realizados, se encuentran versionados mediante Git y alojados en un repositorio público de GitHub:

\begin{center}
\url{https://github.com/Ruipi/tfg}
\end{center}

Los experimentos de entrenamiento se ejecutaron sobre un equipo equipado con una GPU NVIDIA GeForce RTX 5060 Laptop, lo que permitió acelerar significativamente el proceso de entrenamiento del modelo Transformer y realizar múltiples iteraciones experimentales en tiempos razonables.

\begin{table}[H]
\centering

\begin{tabular}{ll}

\toprule
Elemento & Herramienta \\
\midrule

Lenguaje & Python 3.11 \\
Gestión de entorno & Conda (Miniforge) \\
Deep Learning & PyTorch \\
GPU & CUDA \\
Visualización & Matplotlib, Seaborn \\
Análisis de datos & Pandas, NumPy \\
Base de datos & SQLite \\
Desarrollo experimental & Jupyter Notebook \\
Control de versiones & Git + GitHub \\
Editor & Visual Studio Code \\

\bottomrule

\end{tabular}

\caption{Tecnologías empleadas durante el desarrollo del proyecto.}

\end{table}
